%%%%%%%%%%  scan idx 264  |  folio 251  %%%%%%%%%%
\SGAPrintPage{251}%
of a regular immersion and a smooth morphism) is a perfect morphism.

\sgaStableTarget{sga6.e03.notation.4-2}\textbf{Notation 4.2.} Let \(f:X\to Y\) be a locally finite type morphism of schemes. We denote by \(D(f)_{\parf}\), or \(D(X)_{f\text{-}\parf}\), or \(D(X)_{Y\text{-}\parf}\) when no confusion is to be feared (resp. \(D(f)_{\uparf}\), \(D(X)_{f\text{-}\uparf}\), \(D(X)_{Y\text{-}\uparf}\)) the full subcategory of \(D(X)\) formed by the complexes perfect (resp. of finite perfect amplitude) rel. to \(f\). By (I 5.3) and (1.4), \(D(f)_{\parf}\) and \(D(f)_{\uparf}\) are triangulated subcategories of \(D(X)\). When \(X\) is quasi-compact, we have \(D(f)_{\parf}=D(f)_{\uparf}\).

\sgaStableTarget{sga6.e03.prop.4-3}\textbf{Proposition 4.3} (cf. I 5.8). \uline{Let} \(f:X\to Y\) \uline{be a locally finite type morphism of schemes}, \(E\in\Ob D(X)\), \([a,b]\) \uline{an interval of} \(\mathbb Z\), \(x\) \uline{a point of} \(X\), \(y=f(x)\). \uline{The following conditions are equivalent}:
\begin{enumerate}[label=(\roman*)]
\item \(\parfamp_f(E)\subset[a,b]\) \uline{in a neighborhood of} \(x\):
\item \(E\) \uline{is}, \uline{in a neighborhood of} \(x\), \((a-1)\)-\uline{pseudo-coherent rel.} \uline{to} \(f\) \uline{and acyclic outside} \([a,b]\), \uline{and we have} \(\tamp_{\underline{O}_{Y,y}}(E_x)\subset[a,b]\).
\end{enumerate}

\uline{Proof}. The implication (i) \(\Longrightarrow\) (ii) is trivial; let us prove (ii) \(\Longrightarrow\) (i). The question being local in a neighborhood of \(x\), we immediately reduce, by the definitions, to the case where \(f\) is smooth. Then \(E\) is \((a-1)\)-pseudo-coherent in a neighborhood of \(x\); hence, after localizing, we may suppose \(E\) strictly \((a-1)\)-pseudo-coherent (I 2.1) and zero in degree \(>b\), and after truncating (i.e. replacing \(E\) by the complex \((0\to E^a/B^a\to E^{a+1}\to\ldots E^b\to 0)\)), we may suppose that \(E^i\) is free of finite type for \(i>a\), \(E^a\) of finite presentation, and \(E^i=0\) for \(i\notin[a,b]\). By (I 5.1.1), \(E^a_x\) is flat rel. to \(\underline{O}_{Y,y}\), hence (EGA IV 11.3.1) \(E^a\) is \(f\)-flat in a neighborhood of \(x\). Write \(E^a=F\) for brevity; we must show that \(F\) is pseudo-coherent (in a neighborhood of \(x\)). After embedding \(X\) in a standard affine space over \(Y\), we may
%%%%%%%%%%  scan idx 265  |  folio 252  %%%%%%%%%%
\SGAPrintPage{252}%
suppose that \(X=Y[t_1,\dots,t_r]\), with \(Y\) affine and \(f\) the canonical projection. Then \(Y\) is a projective limit of noetherian affine schemes \(Y_\alpha\). Denote by \(f_\alpha:X_\alpha=Y_\alpha[t_1,\dots,t_r]\to Y_\alpha\) the canonical projection. There is (EGA IV 8.5.2 and 11.2.6) an index \(\alpha\) and an \(\underline{O}_{X_\alpha}\)-Module \(F_\alpha\) of finite presentation and \(f_\alpha\)-flat such that \(g_\alpha^*(F_\alpha)\simeq F\), where \(g_\alpha:X\to X_\alpha\) is the canonical morphism. Since \(X_\alpha\) is noetherian, \(F_\alpha\) is coherent, hence pseudo-coherent rel. to \(f_\alpha\). Since, on the other hand, \(F_\alpha\) is \(f_\alpha\)-flat and \(f_\alpha\) is flat, we have \(g_\alpha^*(F_\alpha)\simeq F_\alpha\otimes^{\mathrm L}_{Y_\alpha}\underline{O}_Y\), and it follows from (1.10 a)) that \(F\) is pseudo-coherent rel. to \(Y\), hence pseudo-coherent, q.e.d.

\sgaStableTarget{sga6.e03.cor.4-3-1}\textbf{Corollary 4.3.1.} \uline{Let} \(f:X\to Y\) \uline{be a locally finite type morphism}, \uline{and} \(E\) \uline{an} \(\underline{O}_X\)-\uline{Module}. \uline{For} \(\parfamp_f(E)\subset[0,0]\) \uline{to hold, it is necessary and sufficient that} \(E\) \uline{be}, \uline{rel.} \uline{to} \(f\), \uline{flat and of finite presentation} (1.2.1 c)). \uline{In particular}, \uline{for} \(f\) \uline{to have perfect dimension zero} (4.1), \uline{it is necessary and sufficient that} \(f\) \uline{be flat and locally of finite presentation}.

\uline{Proof}. Apply (4.3 (ii) \(\Longrightarrow\) (i)) with \(a=b=0\).

\sgaStableTarget{sga6.e03.cor.4-3-2}\textbf{Corollary 4.3.2.} \uline{Let} \(f:X\to Y\) \uline{be a locally finite type morphism}, \uline{and} \(E\in\Ob D(X)\) \uline{a complex of} \(\underline{O}_X\)-\uline{Modules flat and of finite presentation} \uline{rel.} \uline{to} \(f\) (1.2.1 c)), \uline{and} \([a,b]\) \uline{an interval of} \(\mathbb Z\). \uline{If} \(E^i=0\) \uline{for} \(i\notin[a,b]\), \uline{we have} \(\parfamp_f(E)\subset[a,b]\). \uline{If} \(E\) \uline{is bounded above in degree}, \(E\) \uline{is pseudo-coherent rel.} \uline{to} \(f\).

\uline{Proof}. A consequence of (4.3.1) and (I 4.15 and 2.11).

\sgaStableTarget{sga6.e03.prop.4-4}\textbf{Proposition 4.4.} \uline{Let there be a commutative triangle of schemes}
\[
\begin{tikzcd}
X \arrow[r,hook,"i"] \arrow[d,"f"'] & Y \arrow[dl,"g"]\\
S &
\end{tikzcd},
\]
%%%%%%%%%%  scan idx 266  |  folio 253  %%%%%%%%%%
\SGAPrintPage{253}%
\uline{where} \(i\) \uline{is a closed immersion and} \(g\) \uline{a smooth morphism}, \uline{and let} \(E\in\Ob D(X)\). \uline{For} \(E\) \uline{to be perfect rel.} \uline{to} \(f\), \uline{it is necessary and sufficient that} \(i_*(E)\) \uline{be perfect} (\uline{in the absolute sense} (I 4.8)). \uline{More precisely}, \uline{suppose} \(g\) \uline{is of relative dimension} \(\leq d\), \uline{and let} \([a,b]\) \uline{be an interval of} \(\mathbb Z\). \uline{Then we have}:
\begin{enumerate}[label=(\roman*)]
\item \(\parfamp(i_*(E))\subset[a,b]\Longrightarrow\parfamp_f(E)\subset[a,b]\);
\item \(\parfamp_f(E)\subset[a,b]\Longrightarrow\parfamp(i_*(E))\subset[a-d,b]\).
\end{enumerate}

\uline{Proof}. An immediate consequence of (3.6).

\sgaStableTarget{sga6.e03.prop.4-5}\textbf{Proposition 4.5.} \uline{Let}
\[
\begin{tikzcd}
X \arrow[rr,"f"] \arrow[dr,"h"'] && Y \arrow[dl,"g"]\\
& S &
\end{tikzcd}
\]
\uline{be a commutative triangle of locally finite type morphisms of schemes}, \uline{and let} \(E\in\Ob D^-(X)\), \(F\in\Ob D^-(Y)\), \([a,b]\), \([c,d]\) \uline{be intervals of} \(\mathbb Z\). \uline{Then we have}:
\[
 \parfamp_f(E)\subset[a,b]\text{ and }\parfamp_g(F)\subset[c,d]
 \Longrightarrow
 \parfamp_h(E\otimes_Y^{\mathrm L}F)
 \subset[a+c,b+d].
\]
\uline{Consequently}, \uline{if} \(E\) \uline{is perfect rel.} \uline{to} \(f\) \uline{and} \(F\) \uline{perfect rel.} \uline{to} \(g\), \(E\otimes_Y^{\mathrm L}F\) \uline{is perfect rel.} \uline{to} \(h\).

\uline{Proof}. Combine (1.11) and (3.4).

\sgaStableTarget{sga6.e03.cor.4-5-1}\textbf{Corollary 4.5.1.} \uline{Under the hypotheses of} (4.5):
\begin{enumerate}[label=\alph*)]
\item \uline{if} \(g\) \uline{has perfect dimension} \(\leq n\), \uline{we have}
\[
 \parfamp_f(E)
 \subset[a,b]
 \Longrightarrow
 \parfamp_h(E)
 \subset[a-n,b];
\]
\uline{consequently}, \uline{if} \(g\) \uline{is perfect}, \(E\) \uline{perfect rel.} \uline{to} \(f\) \uline{implies} \(E\) \uline{perfect rel.} \uline{to} \(h\);
\item \uline{if} \(f\) \uline{has perfect dimension} \(\leq m\), \uline{we have}
\[
 \parfamp_g(F)
 \subset[c,d]
 \Longrightarrow
 \parfamp_h(\mathrm Lf^*(F))
 \subset[c-m,d];
\]
%%%%%%%%%%  scan idx 267  |  folio 254  %%%%%%%%%%
\SGAPrintPage{254}%
\uline{consequently}, \uline{if} \(f\) \uline{is perfect}, \(F\) \uline{perfect rel.} \uline{to} \(g\) \uline{implies} \(\mathrm Lf^*(F)\) \uline{perfect rel.} \uline{to} \(h\).
\end{enumerate}

\sgaStableTarget{sga6.e03.rem.4-5-2}\textbf{Remark 4.5.2.} With the preceding notation, when \(g\) is perfect, \(E\) perfect rel. to \(h\) does not necessarily imply \(E\) perfect rel. to \(f\), as the ever-recurring example of a singular point of a scheme over a field shows. We have the following criterion:

\sgaStableTarget{sga6.e03.prop.4-6}\textbf{Proposition 4.6} (\uline{criterion of perfection by fibers})\SGAAuthorFootnote{1}{Analogous to the ``criterion of flatness by fibers'' (EGA IV 11.3.10).}. \uline{Let}
\[
\begin{tikzcd}
X \arrow[rr,"f"] \arrow[dr,"h"'] && Y \arrow[dl,"g"]\\
& S &
\end{tikzcd}
\]
\uline{be a commutative triangle of morphisms of schemes}. \uline{We suppose} \(f\), \(h\) \uline{locally of finite type}, \uline{and} \(g\) \uline{locally of finite presentation and} flat. \uline{Let} \(x\in X\), \(y=f(x)\), \(s=h(x)=g(y)\). \uline{Finally}, \uline{let} \(E\in\Ob D^-(X)\), \uline{and let} \([a,b]\) \uline{be an interval of} \(\mathbb Z\). \uline{We denote by} \(f_s:X_s\to Y_s\) \uline{the fiber of} \(f\) \uline{at} \(s\), \uline{i.e.} \uline{the morphism obtained from} \(f\) \uline{by the base change} \(\Spec(k(s))\to S\). \uline{Then the following conditions are equivalent}:
\begin{enumerate}[label=(\roman*)]
\item \(\parfamp_h(E)\subset[a,b]\) \uline{in a neighborhood of} \(x\), \uline{and}
\[
 \tamp_{\underline{O}_{Y_s,y}}(E_x\otimes_{\underline{O}_{Y,y}}^{\mathrm L}\underline{O}_{Y_s,y})
 \subset[a,b];
\]
\item \(\parfamp_f(E)\subset[a,b]\) \uline{in a neighborhood of} \(x\).
\end{enumerate}

\uline{Proof}. (ii) \(\Longrightarrow\) (i). We already know (4.3.1) that \(g\) has perfect dimension zero. Thus (4.5.1 a)), condition (ii) implies \(\parfamp_h(E)\subset[a,b]\) in a neighborhood of \(x\). It also implies (by 4.3), \(\tamp_{\underline{O}_{Y,y}}(E_x)\subset[a,b]\), hence (I 5.1 (i)), \(\tamp_{\underline{O}_{Y_s,y}}(E_x\otimes_{\underline{O}_{Y,y}}^{\mathrm L}\underline{O}_{Y_s,y})\subset[a,b]\), which
%%%%%%%%%%  scan idx 268  |  folio 255  %%%%%%%%%%
\SGAPrintPage{255}%
completes the proof of (ii) \(\Longrightarrow\) (i).

(i) \(\Longrightarrow\) (ii). The question being local in a neighborhood of \(x\), we may suppose that \(f\) factors as \(X\overset{i}{\hookrightarrow}X'\overset{f'}{\longrightarrow}Y\), where \(i\) is a closed immersion and \(f'\) a smooth morphism. It is immediate that conditions (i) and (ii) are unchanged when we replace \(E\) by \(i_*(E)\), \(f\) by \(f'\) (and \(h\) by \(gf'\)), so that we may suppose \(f\) smooth. Since \(g\) is perfect (4.3.1), \(h\) is perfect, hence (1.12) \(E\) is pseudo-coherent. After localizing, by (II 2.2.10 b)) and (I 2.10 a)), we may suppose that \(E\) is strictly pseudo-coherent (I 2.1) and zero in degree \(>b\), and that the truncation arrow
\begin{equation}
 E\longrightarrow(0\longrightarrow E^a/B^a\longrightarrow E^{a+1}\longrightarrow\ldots\longrightarrow E^b\longrightarrow0)\tag*{$t)$}
\end{equation}
is a quasi-isomorphism (hypothesis (i) indeed implies that \(E\) is, in a neighborhood of \(x\), acyclic outside \([a,b]\)). We must prove that \(E^a/B^a\) is \(f\)-flat in a neighborhood of \(x\). Denote by \(i:X_s\longrightarrow X\) the canonical arrow. We have:
\begin{equation}
 E_x\otimes_{\underline{O}_{Y,y}}^{\mathrm L}\underline{O}_{Y_s,y}\simeq i^*(E)_x.\notag
\end{equation}
Since \(h\) is flat, hypothesis (i) implies, by (3.5.2), \(\tamp_{h_s}(i^*(E))\subset[a,b]\) in a neighborhood of \(x\). Thus, after localizing, we may suppose that the truncation arrow
\begin{equation}
 i^*(E)\longrightarrow(0\longrightarrow i^*(E^a)/B^a(i^*(E))\longrightarrow\ldots\longrightarrow i^*(E^b)\longrightarrow0)\tag*{$i^*(t)$}
\end{equation}
is a quasi-isomorphism. The hypothesis \(\tamp_{\underline{O}_{Y_s,y}}i^*(E)_x\subset[a,b]\) then implies, by (I 5.1.1), that \(i^*(E^a/B^a)_x\) is a flat \(\underline{O}_{Y_s,y}\)-module. On the other hand, since the arrow \((t)\) is a quasi-isomorphism, the hypothesis \(\tamp_h(E)\subset[a,b]\) implies, again by (I 5.1.1),
